\section{The regularized operator in Dirichlet and Hecke character sectors}
\label{sec:regularized-characters}

This extension concerns actual one-sided closed multipliers and their
smoothed commutators. It carries the construction through the
non-induced imaginary-quadratic example, with the conjugate character
in the reflected component. The actual weighted quotient and cutoff
trace in these sectors have already been supplied by MA-WP.

Fix a number field $K$, a primitive finite-order Hecke character
$\eta$, and a finite set $S$ of finite primes. Extend its ideal
coefficients by zero on ideals not coprime to its conductor, and set
\[
 \begin{aligned}
 L_\eta(s)&=\sum_{\mathfrak a}\eta(\mathfrak a)
                                       (N\mathfrak a)^{-s},\\
 E_{\eta,S}(s)&=\prod_{\mathfrak p\in S}
                       (1-\eta(\mathfrak p)(N\mathfrak p)^{-s}),\\
 A_\eta(s)&=E_{\eta,S}(s)L_\eta(s).
 \end{aligned}
 \tag{MA-RS.40}
\]
The first series is used in $\operatorname{Re}s>1$ and continued
thereafter by the classical Hecke function. Write
$\delta_\eta=1$ for the principal character and zero otherwise,
and $c_\eta=\operatorname{Res}_{s=1}L_\eta(s)$, taking
$c_\eta=0$ when the function is entire. Let $Z_\eta$ be the
nontrivial completed zero divisor, with multiplicities.

\paragraph{MA-RS-005 (actual character commutators and their crossed divisors).}
There is a full-measure subset of $\sigma\in(1/2,1)$ for which
the following assertions hold simultaneously for $\eta$ and
$\overline\eta$. Put $a=\sigma-1/2$ and use the same unitary
$V_a:\mathcal H_a\to L^2(\mathbb R)$ as in (MA-RS.8).
Define the maximal closed multipliers
\[
 Z_{a,\eta,S}=V_a^{-1}M_{A_\eta(\sigma+it)}V_a,
 \qquad
 J_{a,\eta,S}=V_a^{-1}M_{A_\eta(\sigma+it)^{-1}}V_a.
 \tag{MA-RS.41}
\]
For a compact Gevrey test $g$ and a one-sided cutoff $\phi$
with compact Gevrey derivative and endpoints $1,0$,
the actual core composition
\[
        K_{a,\eta,S,\phi}(g)
              =C_g Z_{a,\eta,S}[J_{a,\eta,S},M_\phi]
 \tag{MA-RS.42}
\]
extends to a trace-class operator. With $G=\mathcal Mg$,
\[
\begin{aligned}
 \operatorname{Tr}K_{a,\eta,S,\phi}(g)
  &=-\frac1{2\pi}\int_{\mathbb R}G(a+it)
             \frac{A_\eta'}{A_\eta}(\sigma+it)\,dt\\
  &=\sum_{\substack{\mathfrak p\notin S\\k\ge1}}
       \eta(\mathfrak p)^k\frac{\log N\mathfrak p}
                                    {(N\mathfrak p)^{k/2}}
                         g(k\log N\mathfrak p)
       -\delta_\eta G(1/2)
       +\sum_{\substack{\rho\in Z_\eta\\\operatorname{Re}\rho>\sigma}}
                         m_\rho G(\rho-1/2).
\end{aligned}
\tag{MA-RS.43}
\]
Write the finite prime sum as $P^+_{\eta,\overline S}(g)$ and
put
\[
 B_{\eta,\overline S}(g)
   =P^+_{\eta,\overline S}(g)
                  +P^+_{\overline\eta,\overline S}(g^\vee).
\]
The actual reflected pair on $\mathcal H_a\oplus\mathcal H_{-a}$ is
\[
 \mathbb K_{a,\eta,S,\phi}(g)
   =K_{a,\eta,S,\phi}(g)
                  \oplus R K_{a,\overline\eta,S,\phi}(g^\vee)R.
 \tag{MA-RS.44}
\]
Its trace is
\[
 \operatorname{Tr}\mathbb K_{a,\eta,S,\phi}(g)
  =B_{\eta,\overline S}(g)-\delta_\eta\bigl(G(1/2)+G(-1/2)\bigr)
    +\sum_{\substack{\rho\in Z_\eta\\
                |\operatorname{Re}\rho-1/2|>a}}
                         m_\rho G(\rho-1/2).
 \tag{MA-RS.45}
\]
Every infinite zero sum here is absolutely convergent.

If both divisors instead have a fixed positive horizontal distance
from the chosen line, all these assertions hold for arbitrary
$g\in C_c^\infty$ and $\phi'\in C_c^\infty$, with continuity into
trace norm on fixed compact test supports. Under their Hecke GRH,
this applies at every $1/2<\sigma<1$ and the last sum in (MA-RS.45)
is empty. For Dirichlet characters this uses the stipulated theorem
$A$. For a general non-induced imaginary-quadratic character it uses
its own Hecke zero-location statement; $M$ is sufficient.

\begin{proof}
\emph{Generic inverse estimates and actual domains.}
NT-Z10 and MT-M18 give, for this fixed Hecke function, polynomial
upper bounds on fixed vertical strips, the local zero count
$O(\log(2+|T|))$, and the Euler lower bound
$|L_\eta(2+it)|\ge\zeta_K(2)^{-1}$.
Apply the Borel--Cantelli construction in MA-RS-001 to the union
of the two zero divisors. For almost every $\sigma$, both satisfy
\[
 |\sigma-\operatorname{Re}\rho|
                 \ge c_\sigma(2+|\operatorname{Im}\rho|)^{-2}.
 \tag{MA-RS.46}
\]
At large height the fixed disk centered at $2+it$ contains no
trivial zero or pole of $L_\eta$, and contains only
$O(\log(2+|t|))$ nontrivial zeros. Dividing them out with the
finite Blaschke product and applying Harnack as in (MA-RS.5)
gives the same polynomial lower bound for the zero-free factor.
The product of the bounds in (MA-RS.46) then gives
\[
 |\partial_t^j A_\eta(\sigma+it)^{-1}|
       \le \exp(C_j\log^2(2+|t|))\quad(j\ge0).
 \tag{MA-RS.47}
\]
Cauchy estimates on disks of radius a fixed multiple of
$(2+|t|)^{-2}$ prove the derivative assertion exactly as in
MA-RS-001. All derivatives of $A_\eta(\sigma+it)$ have
polynomial growth. The finite factors $E_{\eta,S}$ and their
inverses have bounded derivatives on a small closed strip about
$\sigma>1/2$, since $|\eta(\mathfrak p)|$ is zero or one.
The compact height interval is covered by nonvanishing on the
line. All statements apply also to $\overline\eta$.

The maximal multiplier operators in (MA-RS.41) are therefore
closed and have the same dense compact-frequency graph core as
in MA-RS.9. Multiplication by $\phi$ sends such a core vector
and its image under $J$ to Fourier tails bounded by
$\exp(-c|t|^{1/q})$, where $q>1$ is a common Gevrey order.
Formula (MA-RS.20) proves this directly, independently of the
arithmetic multiplier. These tails dominate (MA-RS.47), so both
$J M_\phi$ and $Z M_\phi J$ are defined on that core and have
the requisite domain for the next operation.
In weighted Fourier variables the resulting kernel is exactly
\[
 \frac{i}{2\pi}G(a+it)A_\eta(\sigma+it)
  \frac{A_\eta(\sigma+it)^{-1}-A_\eta(\sigma+iu)^{-1}}{t-u}
                 \widehat{\phi'}_+(t-u).
 \tag{MA-RS.48}
\]
The divided difference has subexponential derivative bounds
(MA-RS.47). The two Gevrey factors dominate them in both
variables, as in (MA-RS.22)--(MA-RS.23). Thus this is a Schwartz
kernel, and its Hermite expansion converges in trace norm.
Its diagonal integral, using $\int\phi'=-1$, is the first line
of (MA-RS.43). This proves trace class for the actual composition
before its arithmetic trace is evaluated.

\emph{An explicit ideal-sum regularization.}
The closed operator $Z$ also arises from bounded, explicitly
regularized ideal sums; analytic continuation is not being
substituted for convergence of the unaltered Euler sums.
For $\varepsilon>0$ set
\[
 \begin{aligned}
 L_{\eta,\varepsilon}(s)
     &=\sum_{\mathfrak a}\eta(\mathfrak a)(N\mathfrak a)^{-s}
                                           e^{-\varepsilon N\mathfrak a},\\
 A_{\eta,\varepsilon}(s)
     &=E_{\eta,S}(s)\left(
        L_{\eta,\varepsilon}(s)
             -c_\eta\Gamma(1-s)\varepsilon^{s-1}\right).
 \end{aligned}
 \tag{MA-RS.49}
\]
For fixed $\varepsilon$, these are bounded functions on the
line $\operatorname{Re}s=\sigma$: the ideal series is absolutely
summable there, and $\Gamma(1-\sigma-it)$ is bounded on the
real $t$ line. The finite Euler factor is bounded as well.

Mellin inversion of $e^{-x}$ gives, for $b>\max(0,1-\sigma)$,
\[
 L_{\eta,\varepsilon}(s)
   =\frac1{2\pi i}\int_{b-i\infty}^{b+i\infty}
                  \Gamma(z)L_\eta(s+z)\varepsilon^{-z}\,dz.
 \tag{MA-RS.50}
\]
The interchange with the ideal series is absolute because
$\sigma+b>1$ and $\Gamma(b+iv)$ is integrable.
Fix $0<\theta<\min(1,\sigma)$. Move this contour to
$\operatorname{Re}z=-\theta$.
The crossed poles are $z=0$, with residue $L_\eta(s)$,
and, only for the principal character, $z=1-s$, with residue
$c_\eta\Gamma(1-s)\varepsilon^{s-1}$.
There are no other poles between these lines.
Stirling's exponential vertical decay for $\Gamma$, together
with the fixed-strip polynomial bound for $L_\eta$, makes
the horizontal integrals vanish. The integral on the new
line is $O_T(\varepsilon^\theta)$ uniformly for
$s=\sigma+it$, $|t|\le T$.
Consequently $A_{\eta,\varepsilon}(\sigma+it)$ converges to
$A_\eta(\sigma+it)$ uniformly on every compact frequency
interval. The corresponding bounded multipliers converge on
the compact-frequency core, and the graph closure of the
limit is exactly the maximal $Z$ in (MA-RS.41).
This is an actual regularization with its pole counterterm;
no operator-norm convergence on the whole Hilbert space is claimed.

\emph{The prime and reflected terms.}
On $\operatorname{Re}s=c>1$, the Euler logarithmic derivative is
\[
 -\frac{A_\eta'}{A_\eta}(s)
   =\sum_{\substack{\mathfrak p\notin S\\k\ge1}}
              \eta(\mathfrak p)^k(\log N\mathfrak p)
                                     (N\mathfrak p)^{-ks}.
\]
Absolute convergence and Fourier inversion give
$P^+_{\eta,\overline S}(g)$ for the line integral there.
Shift the line to $\sigma$, using the selected horizontal
heights of MT-M18 for this function. Its minimum-modulus
proof and Stirling bounds give polynomial logarithmic-derivative
bounds on those horizontal segments. The Gevrey Mellin factor
makes their integrals tend to zero.
Each crossed zero has residue $m_\rho$ in $A_\eta'/A_\eta$,
and the pole at one has residue $-\delta_\eta$.
The finite removed Euler factors have their zeros on
$\operatorname{Re}s=0$ and add no crossed term.
The contour orientation is that of MA-RS-003; it gives exactly
(MA-RS.43). The zero count and Mellin decay give absolute
convergence of its residue sum.

The raw character-sector reflection in MA-CC-001 exchanges
$\eta$ and $\overline\eta$. The completed functional equation
identifies the zero divisor of $\overline\eta$ with
$1-Z_\eta$, including multiplicity. Apply (MA-RS.43) to
$\overline\eta$ and $g^\vee$; since
$\mathcal M(g^\vee)(z)=G(-z)$, its right-hand zeros become
exactly the left-hand zeros of $\eta$. Adding gives (MA-RS.45).
A root number is a nonzero constant and introduces no extra
logarithmic-derivative or divisor term.

Finally, with a fixed horizontal gap, every finite Blaschke
factor has a fixed positive lower bound. Since the number of
factors is $O(\log(2+|t|))$, the reciprocal bounds and all their
derivatives are polynomial. The proof of MA-RS-004 then applies
to (MA-RS.48): arbitrary compact smooth tests give a Schwartz
kernel and a continuous trace-class family. Under the stated
Hecke zero-location premise the gap is $\sigma-1/2$ and both
outside sums are empty. This proves all claimed extensions.
\end{proof}

\paragraph{MA-RS-006 (the non-induced cubic sector in the narrow regularization).}
Take the verified character of MA-CH-003:
\[
 K=\mathbb Q(\zeta_3),\qquad
 K(269^{1/3})/K,\qquad \eta=\psi,\qquad
 S=\{\mathfrak p\mid3\cdot269\}.
\]
The test of MA-CH-004 can be chosen of finite Gevrey order.
For that choice it has $g(0)=0$, its local $S$ contribution
is zero, and its full unramified prime contribution is $\log2$.
At every generic width covered by MA-RS-005,
\[
 \operatorname{Tr}\mathbb K_{a,\psi,S,\phi}(g)
     =\log2+
       \sum_{\substack{\rho\in Z_\psi\\
                    |\operatorname{Re}\rho-1/2|>a}}
                       m_\rho\mathcal Mg(\rho-1/2).
 \tag{MA-RS.51}
\]
Under $M$, the same construction works at every $0<a<1/2$
with arbitrary compact smooth versions of this test, and its
trace is exactly $\log2$.

\begin{proof}
The construction in MA-CH-004 uses one bump near norm four,
normalized to have value one there, and subtracts a scalar
multiple of a second bump supported in the interval $(1.4,1.6)$
to cancel the complex archimedean integral. Both bumps can be
chosen compact Gevrey; taking the same scalar combination
preserves that class. The second archimedean integral is nonzero
because its kernel is strictly positive and the bump is chosen
nonnegative and nonzero. All support, value and cancellation
requirements therefore remain exactly those of MA-CH-004.

The sole prime-power contribution in those supports is the
unique prime above two, of norm four and character value one.
Its coefficient is $(\log4)/\sqrt4=\log2$.
The reflected positive-prime sum is zero because $g$ is
supported on positive logarithmic norms. The character is
nonprincipal, so $\delta_\psi=0$. Substitution in (MA-RS.45)
gives (MA-RS.51). Under $M$, all zeros are central, the fixed
gap hypothesis of the smooth extension holds, and the last sum
is empty. This proves the assertion for the actual operator.
It does not use the Dirichlet theorem $A$ to supply a Hecke
premise that $A$ does not contain.
\end{proof}

\paragraph{Character geometry and source locators.}
The raw compact character convention and its inversion under
reflection are those checked in MA-CH-001 and MA-CC-001.
The prime terms are Frobenius traces of the corresponding
rank-one finite-monodromy local system on the open arithmetic
curve where the character is unramified. In the cubic example,
induction to $\mathbb Q$ gives the verified odd rank-two
$S_3$ Artin representation of MA-CH-003.
This is the SGA~1, Expos\'e~V, \S7 finite-\'etale/monodromy
dictionary, and the Artin-motive convention of Deligne,
\emph{Valeurs de fonctions L}, Definition~6.1 and \S\S6.2--6.3.
The actual nuclear trace remains an analytic realization on a
chosen completion, with the cote~1 comparison at inherited
scans~079--080 and~111--112 through BR-TR-01/02.
No positive arithmetic cohomology or new $\ell$-adic weight
statement is inferred from the regularization.

The unconditional Hecke analytic inputs are those of NT-Z10
and MT-M18: Ramakrishnan and Valenza, \emph{Fourier Analysis on
Number Fields}, \S7.4, Theorem~7-19, p.~279, and \S7.7,
Theorem~7-28, p.~289, with the precise strip/count consequences
proved in the inherited modules. The inverse bounds and actual
operator arguments are given above, using the already checked
MA-RS-001--004 kernel proof. The ideal-sum regularization
(MA-RS.49)--(MA-RS.50) uses Mellin inversion of $e^{-x}$:
NIST Digital Library of Mathematical Functions,
(2.5.1)--(2.5.2) and (5.9.1), with $\mu=z=1$ in the latter.
The official formulas were checked on 21 November 2026 at
\url{https://dlmf.nist.gov/2.5.E2} and
\url{https://dlmf.nist.gov/5.9.E1}. The residue computation and
uniform compact-frequency convergence required here are proved
in full, rather than left as an appeal to analytic continuation.
